\section{The shifted crystalline family at every guarded localization}

% BEGIN S80J1
\paragraph{S80J1. The four labelled jump readings for an arbitrary good family.}
Fix a member $D$ with sizes $L$, its original $\mathrm{Facts}$,
and $\mathrm{AdmPar}(L)$. Write $T=(t-1)/2$, $S=S_c$,
$\eta=\eta_{\rm Sz}$ and $b=L.\mathrm{sc}.b$. For every
$1\le j\le m$ and real-place label $\alpha$, set $s=e_{j-1}(\alpha)$.
The upward lists for the cloud, lows and character string are
\[
 C=\{T+s+v_a:0\le a<k-1\},\quad
 A=\{T+s-b_d:0\le d<m\},\quad B=\{0,\ldots,t-1\}.
\]
The downward lists are their images under $z\mapsto2T-z$;
$B$ maps to itself. Write $J=C\cup A\cup B$, $W=C\cup A$,
using the specified orientation, and $H^*=\{-1-z:z\in H\}$ for
the dual of any set of jumps. In these formulas $C,A,B$ have no
repetitions or common terms. Indeed $v_a$ are strictly decreasing in
$[0,S]$, $b_d$ strictly increasing, and $s-b_d$ lies in
$[T+2,\eta]$ for $d<j$ and $[-\eta,-T-2]$ for $d\ge j$.
The lower bound on $s$ is $s>\eta+t+S$ by the three size windows and
$\mathrm{AdmPar}$. Thus the cloud lies above the lows and above $B$,
and the two low ranges lie outside $B$. All these assertions persist
under reflection, with the order reversed.

The $\mathrm{admJ}$ field of the facts, at these parameters, gives
one sign $\mathrm{ep}$ for the datum and the two native reading
identities $\mathrm{AdmDict},\mathrm{AdmRead}$. Take \emph{any}
localization $i$, any good Hecke representation family $R$, any
$w\mid\ell_i$, any coefficient extension of the bridge and any
embedding $\tau:E_w\to L'$ into that coefficient field. The composite
of $\tau|_E$ with the specified complex embedding is a unique host
embedding $g$. The CM half and its conjugate give a real-place label
$\alpha$ over $g$. Use $j=j_i$ and the orientation
$\mathrm{TypeU}(\mathrm{ep},g)$ throughout. The family's de Rham and
level-zero jump identities off the two type primes follow from the first field of $\mathrm{AdmJ}$, and at either type prime
they are its third field. The dictionary gives the first character
a single integral jump $d_\tau=\sigma(-s)-T+c_\tau$, where
$\sigma=1$ or $-1$ according to the orientation. This character
is de Rham at $w$ by the same dictionary and the bridge's
$\mathrm{InducesW}$ clause.

In the full level-zero multiset the entries are $\sigma p+c_\tau$
with $p$ running through $v_a,-b_d,-s+T-c$ for $0\le c<t$.
Subtracting $d_\tau$ gives $T+\sigma(p+s)$. For $\sigma=1$
these are the cloud, the lows and $t-1-c$. For $\sigma=-1$
they are their reflections. Thus $\mathrm{AdmRead}$ identifies, at
every point $x$, the multiplicities in $R_x/\chi_1$ with
$\mathbf1_J(z)$, and those in its Tate-twisted dual with
$\mathbf1_{J^*}(z)$. For any $\pi\in Z_i$ and any native
fraction-field splitting of $R_\pi$ with rank-$r$ complement $\rho$,
the last part of that reading subtracts $\mathbf1_B(z)$ at $\pi$.
It gives $\mathbf1_W(z)$ on $\rho/\chi_1$ and
$\mathbf1_{W^*}(z)$ on $(\rho/\chi_1)^\vee(1)$. The same reading
equates having a jump with nonzero multiplicity for all four
modules. This proves $\mathrm{ConclusionJa}$ with the indicated
single $\mathrm{ep}$ for all families and embeddings.
% END S80J1

% BEGIN S80J2
\paragraph{S80J2. The uniform complete jump bound.}
The signed lists in S80J1 give the remaining six clauses of
$\mathrm{JumpConditions}$. On the upward lists the centered
distances $z-T$ are in $[s,s+S]$ for the cloud,
$T+2\le |z-T|\le\eta$ for the lows, and $|z-T|\le T$ for the
string. The cardinalities of the cloud, lows and entire union are
$k-1,m,n_V$ respectively by disjointness and $n_V=k-1+m+t$.
Upward lows are either in $[T-\eta,-2]$ or $[t+1,T+\eta]$.
The whole cloud lies strictly above $t+1$.
Neither $-1$ nor $t$ belongs to $J$ in either orientation, since
reflection exchanges those two values. Also $0\notin W$ in both
orientations. The numbers of negative and positive lows are
$(m-j,j)$ upward and $(j,m-j)$ downward. Adding in the cloud
gives $m-j$ and $k-1+j$ negative elements of $W$ in the two
orientations.

For the upper bound on $s$, the thick window starts at $b$, the
thin window starts at $w''$, and every top center is at most $b$
by $\mathrm{AdmPar}$. The native upper window allowances give
$s\le b+\eta+1$ for all labels and $s\ge b$ at a thick label.
Thus every $z\in J$ has
\[
 |z|\le b+S+\eta+1+T\le h_J=b+S+\eta+t\le2b,\qquad b\le h_J.
\]
Here $1+T\le t$ and $\eta+t+S<w''\le b$. For a dual cloud element
the two possibilities for its absolute value are $s+v_a+(T+1)$
and $s+v_a-(T+1)$, the latter positive by S80J1.
Thus $s-(T+1)\le |u|\le s+S+(T+1)$. Every dual low satisfies
$-\eta-(T+1)\le u\le\eta-(T+1)$. In either orientation
$0\notin W^*$, and $u=-1-z$ is negative exactly when $z>-1$.
There are consequently $k-1+j$ negative dual values upward and
$m-j$ downward. At a thick label the cloud bounds specialize to
\[
 b-(T+1)\le |u|\le b-(T+1)+(\eta+t+2)+S .
\]
The low bound and the dual string $[-t,-1]$ both give
$|u|\le\eta+t+2$ on the entire dual of the union of lows and string.
This verifies the stated endpoints including the one-unit allowance
in the upper size window.

For the original family form of this theorem, a member at zero
identifies the fixed exponent $\varepsilon$ with a positive scale
exponent. Apply DFD13 using the family's eighth parameter constant
(index $7$, default zero). The parameter tie identifies this constant
with the member's $C_{\rm Lw}$ and the scale variable with $x$ for
$x>\mathrm{Fm}.x_1$. Choose the resulting cutoff once. Every member
beyond it supplied with its $\mathrm{Facts}$ then has $\mathrm{AdmPar}$
and all seven clauses by S80J1 and this computation. This is
$\mathrm{jump\_bound\_of\_facts\_large\_family}$.
% END S80J2

% BEGIN S80J3
\paragraph{S80J3. Crystallinity and residual image for the actual input family.}
Fix $\mathrm{HG}$, a good $R$, $\pi\in Z_i$ and an actual continuous
rank-$r$ complement $\rho$ in the sense of $\mathrm{IsRankRPart}$.
Away from the two type primes, the $\mathrm{g2}$ field of
$\mathrm{Facts}$ gives crystallinity of every $R_x$ at every
$w\mid\ell_i$. The local representation tie identifies this with the
crystalline comparison on the given bridge. Extension of the finite
coefficient field preserves the invariant basis and its dimension,
since taking kernels of linear equations commutes with finite free
scalar extension. We may therefore work over any specified coefficient
extension. The splitting at $\pi$ exhibits $\chi_1$ as a direct
summand since $t\ge1$. The crystalline comparison respects direct
summands, duals and tensor products, using the tied Fontaine
comparison of $\mathrm{Hypotheses}_2$. Hence $\chi_1^{-1}$,
$R_x/\chi_1$, its summand $\rho/\chi_1$ at $\pi$, and the
Tate-twisted dual of that summand are crystalline. The ratio characters
are $\chi_c/\chi_1=\varepsilon_\ell^{-(c-1)}$ for $1\le c\le t$.
The period $t_{\rm cris}$ transforms by $\varepsilon_\ell$ and has
filtration degree one, so $\varepsilon_\ell$ has its unique jump at
$-1$ and the ratio line has its unique jump at $c-1$.

By S80J1 the multiplicities of $R_x/\chi_1$ are independent of $x$
and at most one. By S80J2 it has no jump at $-1$, and neither
$\rho/\chi_1$ nor its Tate-twisted dual has a jump at zero. In either
module a nonzero invariant vector would embed the trivial line.
The filtered de Rham functor on de Rham representations is exact
and strict, so that line would contribute to degree zero at every
embedding, including any one into the chosen coefficient field.
At least one such embedding exists by the bridge coefficient
condition. This gives both vanishing $H^0$ clauses and hence
$\mathrm{GainBoundAt}$ and all of native $\mathrm{admG2}$.

The pointwise $\mathrm{duties}$ in the facts give
$\mathrm{AbsIrrDistinct}$ for every $R_x$ over its fraction field,
with its required decomposition into pairwise inequivalent constituents.
D79G1 gives the residual characteristic polynomial common to all $x$
and the full residual matrix span for the specified $\rho$, and hence
the other two clauses of $\mathrm{admG3}$. The integral matrix entries
of $R_x$ are continuous by goodness, and those of $\rho$ by the
rank-part hypothesis. The intertwining fraction-field isomorphism
of that hypothesis is exactly $\mathrm{SemisimpleDecomposition}$
with the given $t$ ordered character lines. These deductions apply
also to the constant rank-part assignment $x\mapsto\rho$: only its
value at $\pi$ occurs in the assertions.
% END S80J3

% BEGIN S80J4
\paragraph{S80J4. Prescribed labelled characters and unramified twisting.}
For any native case A let $K$ be its unramified local field and $L'$
its prescribed finite coefficient field. Then $K=K_0$, all conjugate
embeddings of $K$ take values in $L'$, and the Wach embeddings
enumerate them bijectively. Denote the embedding at letter $a$ by
$\tau_a$. Given \emph{arbitrary integers} $d_a$, let $\xi$ be the
Lubin--Tate character for the uniformizer $\ell$ of $K$. The
rank-one local comparison used in A72F3 gives the crystalline
character $\tau_a\circ\xi$ with jump $-1$ at $\tau_a$ and $0$ at
all other labels. A72F2 transports this comparison to the native
period rings: the tie intertwines the Galois actions, the inclusion
from crystalline to de Rham periods and the filtration, including the
unramified scalar embeddings specifying the letters. Therefore
\[
                     \theta=\prod_a(\tau_a\circ\xi)^{-d_a}
                        :G_K\longrightarrow O_{L'}^\times
\]
is continuous and crystalline. The values are units since $\xi$
takes values in $O_K^\times$. The tensor product of the corresponding
one-dimensional filtered spaces and inverses has, at $\tau_a$, a
single jump $d_a$ of multiplicity one. Equality of the restrictions
to $K_0=K$ identifies any embedding over letter $a$ with $\tau_a$.

For the second clause of $\mathrm{B3}$, let $\mu:G_K\to (L')^\times$
be continuous and trivial on inertia. Its valuation has compact
image in $\mathbb Z$ and therefore vanishes. Restriction of scalars
on its $O_{L'}$-line gives a $\mathbb Z_\ell$-lattice. In an integral
basis let $M\in\mathrm{GL}_e(\mathbb Z_\ell)$ be the image of an
arithmetic Frobenius, and write $F$ for Witt Frobenius on
$W(\overline{\mathbb F}_\ell)$, $f=[K:\mathbb Q_\ell]$.
There is an invertible matrix $Y$ over this Witt ring satisfying
$M F^f(Y)=Y$. Modulo $\ell$ this is Lang's theorem for
$\mathrm{GL}_e$ over the finite residue field. To lift a solution
modulo $\ell^n$, right multiply an arbitrary lift by $1+\ell^n Z$.
Since $Y^{-1}MF^f(Y)=1\bmod\ell^n$, the new error at order $n$
is the old error plus $F^f(Z)-Z$ modulo $\ell$. This can be made
zero entrywise over the algebraic closure. Completeness gives $Y$.
Its columns give an invariant basis: inertia fixes the unramified
periods, and the Frobenius identity extends by continuity to the
procyclic unramified quotient. The matrices $Y,Y^{-1}$ have entries in the embedded unramified
Witt ring in $B_{\rm cris}$ and map into $B_{\rm dR}^+$. Thus this
change of basis and its inverse preserve every de Rham filtration step. The unramified
line is crystalline with all labelled jumps zero. Tensoring a
crystalline representation with it preserves each labelled graded
dimension and crystallinity. The scalar action specified in
$\mathrm{B3}$ is that tensor product under $v\otimes1\mapsto v$.
Together these statements give all of $\mathrm{adm20}$ for the
native case-A fields, using the specified Wach labels.
% END S80J4

% BEGIN S80J5
\paragraph{S80J5. The coefficient field, inertia, and positive shift at each place.}
Take the common cutoff from S80J2 and G79E5. At a member beyond it
choose $\ell=\ell_i>p_*$ outside the arithmetic exceptional set
and the two type primes. Take any $\mathrm{HG}$, $R,\pi,\rho$
as in S80J3. At each $w\mid\ell$ write $K$ for the completion
presented by its bridge. G79E5 gives $\ell\ge5$ and excludes $q,d$
and the small host primes. The host is unramified at $\ell$ by
its ramification guard, so $K=K_0$. Enlarge the bridge coefficient
field by the normal closure in $\overline{\mathbb Q}_\ell$ of
$K(\mu_\ell)$. This finite compositum $L'$ contains every conjugate
of $K(\mu_\ell)$ and meets exactly the native case-A goodness
conditions at level one. Set $O'=O_{L'}$ and take any generator
$\lambda$ of its maximal ideal. It is irreducible.

The cyclotomic character on inertia at $w$ maps onto
$\mathbb Z_\ell^\times$. Indeed for each $n\ge1$,
$\Phi_{\ell^n}(1+X)$ is Eisenstein over the unramified DVR $O_K$.
The extension by a primitive $\ell^n$th root is thus totally
ramified of degree $\varphi(\ell^n)$. The faithful action on the
root identifies its group with $(\mathbb Z/\ell^n)^\times$.
Absolute inertia surjects onto this finite inertia group. For a
fixed $\ell$-adic unit its prescribed residues define nested
nonempty closed fibers in the compact inertia group. Their
intersection supplies the required element. Transport through
the decomposition-group bridge respects the action on roots of
unity. Since the original exponent $s_0(i)\ge1$, this gives every
$1+\ell^{s_0(i)}z$ and hence $\mathrm{CyclotomicInertiaImage}$.
A place over $\ell$ exists by lying over.

Set $f=[K_0:\mathbb Q_\ell]>0$ and take the actual labels
$\tau_a$ of this field for $a\in\mathbb Z/f$. Apply S80J1 with
$\mathrm{ep}$ to the embeddings $\tau_a$ to obtain their real
places and orientations. G79E5 applied to \emph{this} cycle yields
$h,w',\eta',\mathrm{typ},U$ satisfying $\mathrm{PositionHyps}$
and, with $W_a$ the corresponding oriented union of cloud and lows,
\[
                  \operatorname{count}_{U_a}(z+w'_a)=\mathbf1_{W_a}(z)
                       \qquad (z\in\mathbb Z).
\]
Use S80J4 with $d_a=w'_a$ to get $\theta$ on this same field and
coefficient ring. The position condition gives $w'_a\ge1$.
None of these choices uses the identity of a selected hub point.
% END S80J5

% BEGIN S80J6
\paragraph{S80J6. The full native shifted local condition for every Hecke point.}
With the data of S80J5 form the actual integral local matrices
and lines
\[
 P=\theta(\rho/\chi_1),\qquad D_c=\theta(\chi_c/\chi_1),
 \qquad V_x=\theta(R_x/\chi_1).
\]
All restrictions are through the same bridge and its coefficient
map into $O'$. These are continuous and crystalline by S80J3--S80J4.
At letter $a$, tensoring by $\theta$ moves degree $z$ to
$z+w'_a$ by the filtered tensor rule on the one-dimensional factor.
Thus $P$ has multiplicity $\operatorname{count}_{U_a}(z)$ at $z$,
and $D_c=D_1\varepsilon_\ell^{-(c-1)}$ has a single jump at
$w'_a+c-1$ for $1\le c\le t$. This verifies every field of
$\mathrm{TopHyp}$, since $t\ge1$. At \emph{every} $x$ the
multiplicity of $V_x$ is exactly
\[
       \operatorname{count}_{U_a}(z)
              +\mathbf1_{[w'_a,w'_a+t-1]}(z)
                =\mathrm{jTil}(U,w',t)(a,z).
\]
Here the disjoint union with the unshifted string is S80J1,
including in the downward orientation.

For $g\in G_K$, the residual polynomial equality of S80J3 passes
through the residue map to $O'/(\lambda)$. Indeed the bridge
identifies the original integer ring with that of its coefficient
field and extension takes positive valuation to positive valuation.
Twist each matrix by the same unit $\theta(g)/\chi_1(g)$.
The coefficient of $X^{n_V-j}$ in a monic characteristic polynomial
of degree $n_V$ scales by its $j$th power. Applying the same root
scaling to the product on the right gives, for every $x,g$,
\[
 \det(X-V_x(g))\equiv
 \det(X-P(g))\prod_{c=1}^t(X-D_c(g))\pmod\lambda .
\]
Every coefficient of the difference is integral and has the form
$\lambda y$ with $y\in O'$, as required also after embedding in
$L'$. This proves $\mathrm{IsFull}$ for every $V_x$. Together with
$\mathrm{TopHyp}$ this is precisely $\mathrm{SetC}$ of the original
$\mathrm{exists\_exactLocalConditionData\_at\_place}$. More generally
the same last construction works at any place with the literal local
input hypotheses of that theorem, using its supplied position
witnesses and jump reading. In that formulation disjointness of the
shifted string from $U_a$ also follows directly from position. A
shared point would belong to the low window since the whole string
is in that window, yet its distance from itself would have to exceed
the nonnegative carry plus two. Under the original family and prime
guards, S80J1--S80J5 furnish those inputs on every guarded
$(R,\pi,\rho)$. The further boundedness conclusion for cut cocycles
uses a separate local cohomology criterion.
% END S80J6
